\par\noindent\textbf{(1)} 不可压无旋给 \(\nabla^2\Phi=0\)，轴对称 \(l=1\) 一般解 \(\Phi=(Ar+B/r^2)\cos\theta\)。无穷远 \(A=U\)，球面 \(\partial_r\Phi|_{r=a}=0\) 给 \(B=Ua^3/2\)。故 \(\Phi=U(r+a^3/(2r^2))\cos\theta\)。
\par\noindent\textbf{(2)} 径向速度 \(v_r=U(1-a^3/r^3)\cos\theta\)，角向速度 \(v_\theta=-U(1+a^3/(2r^3))\sin\theta\)。球面 \(r=a\) 上 \(v_r=0\)、\(v_\theta=-3U\sin\theta/2\)，在上下游两极 \(\theta=0,\pi\) 形成理想驻点。
\par\noindent\textbf{(3)} 沿定常无旋流 \(p+\rho v^2/2=p_\infty+\rho U^2/2\)。因此 \(p(a,\theta)=p_\infty+\tfrac12\rho U^2[1-(9/4)\sin^2\theta]\)。它关于 \(\theta\to\pi-\theta\) 对称，不能产生前后压差阻力。
\par\noindent\textbf{(4)} 球的流向压力力 \(F_z=-\int_0^\pi p(a,\theta)\cos\theta\,2\pi a^2\sin\theta d\theta=0\)，因为含 \(1\) 与 \(\sin^2\theta\) 的被积函数积分均为零。此D’Alembert零阻力依赖无黏无分离的理想势流；真实Reynolds数下边界层与尾涡产生阻力。
